% 18 November 2026. Prepared for Edward Baker with OpenAI Codex assistance.
% Actual model/effort unavailable; inherited helper settings, no attestation.
\section{Joint forcing dynamics and deflated logarithmic derivatives}
\label{sec:joint}

The direct forcing bound in Section~\ref{sec:bounds} estimates each moment
separately. The hierarchy permits a more structured question: what is the
exact evolution of their forcing combination, and what does its coupling
to the gap actually control? We first keep that combination intact, then
identify an equivalent representation in derivatives of the pair-deflated
field. The identities and the two scoped tests below develop
\cite{BakerCM13}; they retain the strong-collar tube
\eqref{eq:domain} throughout. They supply no additional canonical sign
estimate.

\subsection{An exact derivative formula for the forcing combination}

Use the fixed normalized variables
\begin{equation}
 s_n=\ell^{2n}S_n=Z(n),\qquad
 v_n=\ell^{2n-1}T_n,\qquad
 m=s_1,\quad \delta=D/\ell^2,\quad
 \sigma=(t-t_0)/\ell^2.
 \label{joint:normalization}
\end{equation}
Here $s_n$ is a moment, whereas the unindexed $s$ in $z_\pm=c\pm s$
remains a local square root of $D$. In this subsection and the next,
primes denote $\sigma$ derivatives. The forcing and gap equations are
\begin{equation}
 F=12s_2+16\delta s_3,\qquad m'=2m^2-F,\qquad
 \delta'=2-4m\delta.
 \label{joint:system}
\end{equation}

\begin{proposition}[Joint forcing remainder]\label{joint:forcing-prop}
On the complete tube, including mergers,
\begin{align}
 F'-8mF={}&\mathcal R,\label{joint:forcing-law}\\
 \mathcal R={}&48v_2^2+192\delta v_2v_3+144\delta s_2^2
       +192\delta^2s_2s_3\notag\\
       &\quad-352s_3-1440\delta s_4-1152\delta^2s_5.
       \label{joint:remainder}
\end{align}
Consequently, with the multiplier $W$ of \eqref{eq:transport-cost},
\begin{equation}
 F(y)=W(x,y)^{-2}F(x)
       +\int_x^y W(u,y)^{-2}\mathcal R(u)\dd u.
 \label{joint:forcing-integral}
\end{equation}
\end{proposition}

\begin{proof}
The second and third even rows of Theorem~\ref{thm:hierarchy}, after the
constant normalization \eqref{joint:normalization}, read
\begin{align*}
 s_2'&=8ms_2-32s_3+4\delta s_2^2-40\delta s_4+4v_2^2,\\
 s_3'&=12ms_3+6s_2^2+12v_2v_3+12\delta s_2s_3
                   -60s_4-72\delta s_5.
\end{align*}
Apply the product rule to $F$ and substitute $\delta'=2-4m\delta$.
The terms involving $m$ are
$96ms_2+(-64+192)m\delta s_3=8mF$.
The coefficients of $s_3$, $\delta s_2^2$ and $\delta s_4$ are,
respectively, $-384+32=-352$, $48+96=144$ and $-480-960=-1440$.
The other terms give \eqref{joint:remainder}. Only five fixed even orders
and two fixed odd orders occur; their analytic contour expressions and
complete convergence were proved in Theorem~\ref{thm:hierarchy}.
Thus no growing-order or differentiated sharp-cutoff passage is used.
Finally the integrating factor for \eqref{joint:forcing-law} is
$\exp(-8\int_x^\sigma m)=W(x,\sigma)^2$. Integrating and dividing by
$W(x,y)^2$ gives \eqref{joint:forcing-integral}.
\end{proof}

This exact combination makes the missing estimate visible. If
$u_j=w_j/\ell$ and $\mu_j=\ell^2/Q_j$, then
$|u_j||\mu_j|^{1/2}\le\sqrt{q_0}$ and $|\mu_j|\le\gamma$.
Therefore \eqref{eq:integer-bound} gives
\begin{equation}
 |v_n|\le\sqrt{q_0}\,\gamma^{n-3/2}\mathcal A_1
          \le\sqrt{q_0}\,\gamma^{n-3/2}B_1L,
 \qquad n\ge2.
 \label{joint:odd-bound}
\end{equation}
At fixed collar parameters these separate bounds produce an $O(L^2)$
allowance for $\mathcal R$, not a smaller forcing scale. This is a
statement about that allowance, not a lower bound on the actual remainder.
Nor is there a useful sign hidden in the first two terms: their completion
is $48(v_2+2\delta v_3)^2-192\delta^2v_3^2$. A signed use of
\eqref{joint:forcing-integral} must also control its entry value and
weights; assuming the desired transport bound for those weights would
not establish it.

\begin{remark}[Exact local sign test]\label{joint:quartic-test}
For a finite polynomial test only, take $\ell=1$ and initialize
$P(0,z)=z^2(z^2+b^2)$, $b>0$. Its exact heat evolution is
$P(0,z)-tP_{zz}(0,z)+t^2P_{zzzz}(0,z)/2$. A small time neighborhood
isolates the two occurrences at the origin from the exterior roots
$\pm\ii b$. At the initial double,
\[
 m=-2/b^2,\quad s_2=2/b^4,\quad s_3=-2/b^6,\quad v_2=0,
 \qquad F=24/b^4,\quad \mathcal R=704/b^6.
\]
Initializing instead $z^2(z^2-b^2)$ changes the exterior roots to
$\pm b$ and gives $\mathcal R=-704/b^6$. Thus the finite hierarchy,
conjugation and local collar alone do not force $\mathcal R\le0$.
These polynomials lack the unchanged theta source, its positive value at
the origin and its growing complete population. They are local tests of
the algebra, not counterexamples to a canonical estimate.
\end{remark}

\subsection{The regular gap-energy identity and its weight}

A second way to retain the coupling is to eliminate $m$ from
\eqref{joint:system}. Division by $\delta$ would introduce an artificial
singularity at a merger. The following identity avoids it.

\begin{proposition}[Signed weighted gap energy]\label{joint:energy-prop}
On every closed scaled subinterval $[x,y]$ of the tube,
\begin{align}
 2\delta\delta''-(\delta')^2+4&=8F\delta^2,
       \label{joint:energy-point}\\
 8\int_x^yF\delta^2\dd u
 &=2[\delta\delta']_x^y-3\int_x^y(\delta')^2\dd u+4(y-x)
       \label{joint:energy-integral}\\
 &\le2[\delta\delta']_x^y+4(y-x).\notag
\end{align}
These formulas hold through $\delta=0$.
\end{proposition}

\begin{proof}
Differentiating the last equation of \eqref{joint:system} and substituting
the other one gives
$\delta''=8m^2\delta-8m+4F\delta$.
Multiplication by $2\delta$ and subtraction of
$(\delta')^2=(2-4m\delta)^2$ proves \eqref{joint:energy-point}.
All terms are regular, including at a zero of $\delta$.
Integrate and use
$\int_x^y\delta\delta''=[\delta\delta']_x^y-
\int_x^y(\delta')^2$ to obtain \eqref{joint:energy-integral}.
Reality of $\delta'$ proves the inequality.
\end{proof}

The boundary value is $\delta\delta'=2\delta-4m\delta^2$ and hence
still contains the interaction. The weight vanishes at mergers.
Even on an interval where it is bounded below, an upper bound for a
weighted \emph{signed} integral cannot be divided by the minimum weight
to bound the unweighted integral without controlling its negative part.
The exact transport cost in \eqref{eq:transport-cost} involves a different
weight.

\begin{proposition}[Limitation in the scalar consumer class]
\label{joint:scalar-screen}
The weighted allowance \eqref{joint:energy-integral}, together with the
two scalar equations \eqref{joint:system}, does not bound $W$ or its
inhomogeneous cost uniformly, even when $|\delta|\le d$ for fixed $d>0$.
This conclusion concerns only scalar solutions of those two equations.
\end{proposition}

\begin{proof}
For $h\ge1/(2d)$ set $m=-h$, $F=2h^2$ and $\delta=-1/(2h)$.
Both equations hold identically. The weighted integral is
$\int_x^yF\delta^2\dd u=(y-x)/2$, and
\eqref{joint:energy-integral} is an equality independent of $h$.
In contrast,
\begin{equation}
 W(x,y)=\e^{4h(y-x)},\qquad
 2\int_x^yW(u,y)\dd u=\frac{\e^{4h(y-x)}-1}{2h},
 \label{joint:scalar-cost}
\end{equation}
which are unbounded as $h\to\infty$ at fixed $y-x>0$.
\end{proof}

No divisor or full-hierarchy realization of this scalar family is asserted.
It therefore does not exclude an estimate using their additional structure.
For comparison with the baseline $C_FL$ allowance, taking
$h=\sqrt{C_FL/2}$ retains that forcing scale for large $L$ and the same
exponential transport scale. The linearization
$q=\exp(-2\int m)$, $q''=2Fq$ and $W(x,y)=[q(y)/q(x)]^2$ is equivalent
to the scalar system and adds no signed information.

\subsection{Deflation and an absolutely convergent derivative expansion}

The same forcing also has a representation using only local derivatives
of an entire function with the target pair removed. In the moving
coordinate $w$, define
\begin{equation}
 G(t,w)=\frac{H(t,c(t)+w)}{w^2-D(t)},\qquad
 U_n(t)=\sum_j w_j(t)^{-n},\quad n\ge2.
 \label{joint:deflated-field}
\end{equation}
The quotient removes both target occurrences, also at $D=0$, and is
entire in $w$. It is zero-free on $|w|<r_0$. Thus a holomorphic logarithm
exists there; its derivatives of positive order do not depend on the
chosen additive logarithm constant.

\begin{proposition}[Deflated derivative transfer]\label{joint:deflation-prop}
For $n\ge2$,
\begin{equation}
 U_n=-\frac{\partial_w^n\log G(t,0)}{(n-1)!}.
 \label{joint:U-derivative}
\end{equation}
For each fixed integer $n\ge1$,
\begin{align}
 S_n&=\sum_{k\ge0}\binom{n+k-1}{k}D^kU_{2n+2k},
       \label{joint:S-expansion}\\
 F&=4\ell^4\sum_{k\ge0}(k+1)(2k+3)D^kU_{2k+4}.
       \label{joint:F-expansion}
\end{align}
The series are absolutely convergent, uniformly on the established real
time tube. Truncation of \eqref{joint:F-expansion} after $k=K$ has error
at most
\begin{equation}
 4\ell^4\left(\sum_j|w_j|^{-4}\right)
       \sum_{k>K}(k+1)(2k+3)q^k,
 \qquad q=|D|/r_0^2\le1/16.
 \label{joint:deflation-tail}
\end{equation}
In particular, at a merger,
\begin{equation}
 F=12\ell^4U_4=-2\ell^4\partial_w^4\log G(t,0).
 \label{joint:merger-derivative}
\end{equation}
\end{proposition}

\begin{proof}
The deflated entire function has order at most one and precisely the
exterior zeros $w_j$. Its genus-one logarithmic product has a possible
affine prefactor and linear regularizers. Their derivatives vanish after
two differentiations. Differentiating the remaining factors at $w=0$
gives $-(n-1)!w_j^{-n}$, proving \eqref{joint:U-derivative}.
The sum and its derivatives converge on smaller disks by the complete
inverse-square tail and the exterior clearance.

For each occurrence, the collar gives $|D/w_j^2|\le1/16$. Apply the
binomial series to
$Q_j^{-n}=w_j^{-2n}(1-D/w_j^2)^{-n}$.
The absolute double sum is bounded by
\[
 \sum_j|w_j|^{-2n}
       \sum_{k\ge0}\binom{n+k-1}{k}16^{-k}
   =(1-1/16)^{-n}\sum_j|w_j|^{-2n}.
\]
The last sum is bounded uniformly on the fixed tube by the collar and
unit count. This proves \eqref{joint:S-expansion} with the stated
complete and time-uniform passage. Since
$F=\ell^4(12S_2+16DS_3)$, the coefficient of $D^kU_{2k+4}$ is
$12(k+1)+8k(k+1)=4(k+1)(2k+3)$, including $k=0$.
This proves \eqref{joint:F-expansion}. Taking absolute values only in
its omitted terms gives \eqref{joint:deflation-tail}. At $D=0$ retain
$k=0$ and use $3!=6$ in \eqref{joint:U-derivative}.
\end{proof}

\subsection{The fixed-source superzeta correspondence}

For clarity about the relation to classical superzeta functions, denote
zeros of $\xi$ by $\varrho$ and write the first-kind integer sum as
\[
 \mathfrak Z(n\mid\tau)=\sum_{\varrho}
              (1/2+\tau-\varrho)^{-n},\qquad n\ge2.
\]
This is the zero-based notation of \cite{Voros2014SuperzetaReview}; an
irrelevant constant normalization of the completed function does not
change it. If the independent pair tube contains $t=0$ and $D\ne0$, the
normalization \eqref{eq:source} gives
\begin{equation}
 U_n(0)=(\ii/2)^n\mathfrak Z(n\mid-\ii c(0)/2)
                     -s(0)^{-n}-(-s(0))^{-n}.
 \label{joint:superzeta-transfer}
\end{equation}
Indeed, for a zero $z_j=c+w_j$ of $H(0,\cdot)$ and its corresponding
$\varrho=(1-\ii z_j)/2$, one has
$(1-\ii c)/2-\varrho=\ii w_j/2$. Summing over the complete divisor
and removing the two target terms proves the formula. For even $n=2r$
the removed term is $2D^{-r}$. These separate pieces become singular at
a merger; the deflated expression \eqref{joint:U-derivative} is the
regular combination. Separate absolute bounds for the pieces would lose
that cancellation.

The auxiliary translation $\tau$ in a fixed-source superzeta sum is not
heat time. To see the distinction directly, put
$\Xi_t(\zeta)=8H(t,\ii(2\zeta-1))$. The chain rule and
\eqref{eq:heat} give
\begin{equation}
 \Xi_0=\xi,\qquad \partial_t\Xi_t=\tfrac14\partial_\zeta^2\Xi_t.
 \label{joint:Xi-heat}
\end{equation}
Defining a superzeta for this heat-deformed field returns the same local
derivatives in \eqref{joint:U-derivative}. The fixed-$\xi$ continuation
formulas of \cite{Voros2014SuperzetaReview} do not themselves provide a
signed estimate for these moving-center, pair-deflated derivatives on a
heat interval. Moreover, the integer forcing
\eqref{joint:F-expansion} already lies in the domain of absolute
convergence; continuation to $\Re p\le1/2$ is unnecessary for this
identity.

Together, \eqref{joint:forcing-integral}, \eqref{joint:energy-integral}
and \eqref{joint:F-expansion} distinguish three exact reorganizations
of the same quantitative problem. A useful further estimate must control
a signed remainder with its weights and entry value, a nondegenerate
quantity that bounds the transport cost, or the joint deflated derivatives
on the established tube. None of those estimates follows merely from the
representations, and direct residues have access to the same information.
